\section{Basic Deployment Pipeline - Modeling with Petri Nets}
\subsection{a)} 

I asked ChatGPT to specify three types of tests for the game described above.

The specific prompt I used was:

\begin{adjustwidth}{2em}{0pt}
{
Suppose I am developing an action-adventure game such that:
a) players explore a post-apocalyptic version of McMaster University; b) up to eight players 
may be accosted by zombies, mutated professors, angry Deans, killer squirrels, and more; c) 
gameplay can be both competitive and cooperative; and d) the game will be released on the 
PS5, Xbox Series X, Nintendo Switch and Sega Genesis (which is 16-bit and hasn't been 
manufactured since 1997).

Specify three types of tests for such a game.
}
\end{adjustwidth}

The response I received was as follows:

\begin{adjustwidth}{2em}{0pt}
{

For a game like this — especially one that spans \textbf{modern consoles and a 16-bit relic like the Sega Genesis} — testing needs to cover functionality, performance, and platform constraints very carefully.

Here are \textbf{three distinct and appropriate types of testing}:

1. Functional Testing (Gameplay \& Mechanics)

\textbf{Purpose:} Ensure the game works according to design.

This would verify:

\begin{itemize}
    \item Players can explore the post-apocalyptic McMaster campus without falling through geometry
    \item Zombies, mutated professors, angry Deans, and killer squirrels spawn correctly
    \item Combat, health, inventory, and abilities function as intended
    \item Competitive and cooperative modes behave correctly
    \item Up to eight players connect and interact properly
\end{itemize}

}
\end{adjustwidth}

\section{Safety and other useful properties - Formalizing requirements }

\section{Critical updates - Working with priorities}

\section{Formalizing Requirements}

\section{Reflection}